\documentclass[11pt]{article}
\usepackage[utf8]{inputenc}
\usepackage[T1]{fontenc}
\usepackage[a4paper,margin=1in]{geometry}
\usepackage{lmodern}
\usepackage{microtype}
\usepackage{amsmath,amssymb,mathtools}
\usepackage{enumitem}
\usepackage{longtable,booktabs,array}
\usepackage[hidelinks]{hyperref}
\setlength{\parindent}{0pt}
\setlength{\parskip}{0.55em}
\newcommand{\K}{K}
\newcommand{\Gr}{\operatorname{Gr}}
\newcommand{\Fil}{\operatorname{Fil}}
\newcommand{\cl}{\operatorname{cl}}
\newcommand{\ch}{\operatorname{ch}}
\newcommand{\Todd}{\operatorname{Todd}}
\newcommand{\Parf}{\operatorname{Parf}}
\newcommand{\Inv}{\operatorname{Inv}}
\newcommand{\Pic}{\operatorname{Pic}}
\newcommand{\Spec}{\operatorname{Spec}}
\newcommand{\codim}{\operatorname{codim}}
\newcommand{\dimx}{\operatorname{dim}}
\newcommand{\RRR}{\mathrm{RRR}}

\begin{document}
\begin{center}
{\Large SGA 6 -- Expos\'e 0}\par
{\large Outline of a Programme for a Theory of Intersections on General Schemes}\par
{\large continuation of Sections 3--7 and bibliography}
\end{center}

If $f$ is a locally complete intersection morphism, one introduces a natural notion of ``virtual relative dimension'' $d(f)$ by the formula
\[
d(f)(x)=\dim_x\bigl(f^{-1}(f(x))\bigr)-\codim_x(X,X').
\]
It is an integer-valued function, locally constant at the point $x\in X$, hence constant if $X$ is connected; moreover it does not depend on the choice of the factorization (2.1). When this virtual relative dimension $d$ is constant, when $f$ is proper, and when $X$ and $Y$ admit ample invertible modules, one can prove (Expos\'e VIII), although this is far from a trivial result, that
\begin{equation}\tag{3.4}
f_*\bigl(\Fil^i K^\bullet(X)_{\mathbb Q}\bigr)\subset \Fil^{i-d}K^\bullet(Y)_{\mathbb Q},
\end{equation}
whence a homomorphism (3.3) of degree $-d$.

\subsubsection*{3.3}
We have therefore assembled all the structural elements needed to give a meaning to formula (1.2) in the case of a proper locally complete intersection morphism of noetherian schemes $X,Y$, both admitting ample invertible modules. With a little additional care, this formulation in fact retains a meaning independently of any noetherian hypothesis.

In this general form, this Riemann--Roch formula will be proved in Expos\'e VIII. The reader will observe that the proof given in this seminar combines elements of the two proofs which appear in [RRR] and which were discussed in the Introduction. The need to resort to the methods of the first of these proofs, involving more $K$-formalism than the second, followed in Borel--Serre, is due to the need for a proof of the inclusion (3.4), a question which did not arise in the more restricted framework of [RRR] and Borel--Serre [2], where one could work with the Chow ring or the ring $\Gr^\bullet_{\mathrm{top}}$. From the second proof, we retain the introduction of the scheme obtained by blowing up $Y$ along $X$ in the case where $f$ is an immersion. This technical artifice will no doubt disappear, and the proof in the case of an immersion will extend to the case where one no longer postulates the existence of an ample invertible module on $Y$, once the question raised in Expos\'e XIV, no. 1, concerning the introduction of exterior-power operations in the derived category $D^-(Y)$ has been resolved.

\subsection*{4. Removing the ample invertible module hypothesis}
We must finally give a few indications of how one can still give a meaning to (1.2) in the absence of a hypothesis asserting the existence of ample invertible modules on $X$ and $Y$.

\subsubsection*{4.1}
First note that for an arbitrary scheme $X$, or more generally an arbitrary locally ringed topos, it seems beyond doubt that the ``right'' invariant which should replace the ring $K(X)$ of [RRR] is not the ring $K^\bullet(X)$ considered above, defined in terms of locally free modules on $X$, but rather the group $K(\Parf(X))$ considered in 3.1, defined in terms of perfect complexes on $X$. We shall denote this group here by $K^\bullet(X)$, and distinguish it from the ``naive'' group of 3.1, denoted $K^\bullet_{\mathrm{naive}}(X)$. This point of view is forced by the need to define a direct-image homomorphism (3.1) for a proper morphism
\[
f:X\longrightarrow Y
\]
of finite Tor-dimension, say between noetherian schemes; the non-noetherian case requires somewhat stronger local finiteness properties on $f$, cf. Expos\'e III. With the definition of $K^\bullet(X)$ now adopted, this homomorphism $f_*$ is indeed defined in a trivial way by formula (3.2), where now $F$ is an arbitrary perfect complex on $X$, which implies that $Rf_*(F)$ is also perfect.

\subsubsection*{4.2}
This new definition of $K^\bullet(X)$ raises, however, a new problem in homological algebra: that of defining a $\lambda$-structure on $K^\bullet(X)$ compatible with its natural ring structure coming from tensor product in the derived category. If $L$ and $L'$ are two perfect complexes, then so is their total tensor product $L\otimes^{\mathbf L}L'$. More precisely, one would need to define ``exterior-power'' operations in the derived category $D^-(X)$, transforming perfect complexes into perfect complexes. This question has not yet been clarified at present, and, as was said in the Introduction, it constitutes the most serious gap in the seminar. Once this question is solved, and proceeding as in 2.3, one obtains a natural filtration on $K^\bullet(X)$, allowing one to form the associated graded ring $\Gr^\bullet(X)$, which will play the role of the Chow ring, and to define a theory of Chern classes
\[
c_i:K^\bullet(X)\longrightarrow \Gr^i(X).
\]

\subsubsection*{4.3}
To give a meaning to formula (1.2) for a proper locally complete intersection morphism
\[
f:X\longrightarrow Y
\]
of noetherian schemes, one must still include in the statement of this formula the nontrivial inclusions (3.4), where $d$ is the virtual relative dimension, in order to be able to define (3.3); and finally one must define the virtual relative tangent class
\[
T_f\in K^\bullet(X).
\]
When $f$ admits a factorization (2.1) through an immersion into a relative smooth scheme, the definition of $T_f$ presents no new difficulty. It can moreover be reformulated in that case, in a way more consonant with the spirit of derived categories, by considering the well-known canonical homomorphism
\begin{equation}\tag{4.1}
N_{X/X'}\longrightarrow \Omega^1_{X'/Y}\otimes_{\mathcal O_{X'}}\mathcal O_X
\end{equation}
induced by the exterior differential, and the chain complex $L_{X/Y}$ defined from (4.1) by ``extending by zeros'', with the complex equal in degree $1$, respectively $0$, to the source, respectively the target, of the homomorphism (4.1). Formula (2.2) then takes the more satisfactory form
\begin{equation}\tag{4.2}
T_f=\cl\bigl(L_{X/Y}^{\vee}\bigr),
\end{equation}
a formula which makes sense in $K^\bullet(X)$ when $f$ is a locally complete intersection morphism, since in that case the complex $L_{X/Y}$ is manifestly perfect. The assertion that $T_f$ is invariant relative to the chosen factorization of $f$ can then be made precise as an invariance assertion for the complex $L_{X/Y}$: up to canonical isomorphism in the derived category $D(X)$, this complex is likewise independent of the chosen factorization (Expos\'e VIII).

\subsubsection*{4.4}
When $f:X\to Y$ is an arbitrary morphism of schemes, one can still construct a canonical object $L_{X/Y}$ of the derived category $D(X)$, defined up to unique isomorphism, which on each affine open of $X$ lying above an affine open of $Y$ is isomorphic to the complex constructed according to the principle explained in the preceding paragraph. In the general case, however, ``smooth'' should be replaced by ``formally smooth'', so that every algebra $B$ over a ring $A$ occurs as a quotient of a formally smooth algebra, for example a polynomial algebra. For such a definition it is not possible to proceed by simple gluing from the affine constructions, since the objects of the derived category $D(X)$ are, essentially, not glueable.

Here is a general construction of $L_{X/Y}$, valid more generally for every morphism
\[
f:X\longrightarrow Y
\]
of commutatively ringed topoi. Let $B=\mathcal O_X$ and $A=f^{-1}\mathcal O_Y$, so that $A$ is a ring on the space $X$ and $B$ is an $A$-algebra. Let $T\to B$ be a homomorphism of sheaves of sets which generates $B$ as an $A$-algebra, i.e. such that the corresponding homomorphism of $A$-algebras
\[
C=A[T]\longrightarrow B
\]
is an epimorphism of sheaves of sets; here $A[T]$ denotes the commutative $A$-algebra freely generated by $T$, that is, the sheaf associated with the presheaf
\[
U\longmapsto A(U)[T(U)].
\]
Let $J$ be the kernel of $C\to B$, and consider, as in 4.3, the chain complex of length $1$ defined by the natural homomorphism
\[
J/J^2\longrightarrow \Omega^1_{C/A}\otimes_C B.
\]
It is not difficult to see that, up to canonical isomorphism in the category $D(X)$, the complex so constructed is independent of the choice of $T$ and of $T\to B$, so that it deserves the designation of the relative cotangent complex $L_{X/Y}$ of $X$ over $Y$. One also verifies that, on every affine open of $X$ over an affine open of $Y$, it is given, up to canonical isomorphism, by the procedure of 4.3.

In the case where $f$ is locally complete intersection, one sees moreover that the result is a perfect complex, which makes it possible again to define an element $T_f$ of $K^\bullet(X)$ by formula (4.2). Here again, in general, it would not have been possible to define $T_f$ as an element of $K^\bullet_{\mathrm{naive}}(X)$, and nothing permits one to suppose that $T_f$ is globally isomorphic in $D(X)$ to a bounded complex with locally free components of finite type. This is therefore an additional justification for the point of view advanced in 4.1 in favour of the ``sophisticated'' $K^\bullet(X)$, in preference to $K^\bullet_{\mathrm{naive}}(X)$.

\subsubsection*{4.5}
Subject to the question raised in 4.2, we have therefore again assembled all the structural elements needed to give a meaning to (1.2) for a proper locally complete intersection morphism of arbitrary schemes; the noetherian hypothesis implicitly made in the preceding sections can in fact be eliminated without any essential difficulty. Unfortunately, this does not mean that, even subject to this reservation, the formula in question is established, even in the case where $Y$ is the spectrum of an algebraically closed field of characteristic $p>0$ and $X$ is, say, a proper smooth algebraic scheme of dimension $3$. See the comments in Expos\'e XIV, no. 2.

\subsection*{5. Analytic variant}
In the form presented in no. 4, the Riemann--Roch formula (1.2) also has a meaning for a proper locally complete intersection morphism of complex analytic spaces, or of rigid analytic spaces in the sense of Tate. But of course the proof given in [RRR] applies to this case only when one assumes in addition that the morphism $f$ is projective; and, as in the algebraic case, it seems that an essentially new idea is needed to treat the general case. The formula we have in view here is a formula with values in an ``analytic Chow ring'' $\Gr^\bullet(X)$, or rather $\Gr^\bullet(X)_{\mathbb Q}$, constructed according to the principle explained in 4.2, the difficulty encountered there disappearing, moreover, because one is in characteristic zero, as is noted in Expos\'e XIV, no. 1.

When $X$ and $Y$ are nonsingular, it is possible to give another version of the Riemann--Roch formula by proceeding as follows. First let $X$ be an arbitrary complex analytic space, and let $X^{\mathrm{top}}$ be the same space equipped with the sheaf of rings of continuous complex-valued functions. We therefore have a canonical morphism of ringed spaces
\[
X^{\mathrm{top}}\longrightarrow X,
\]
which, by the evident contravariant nature of the functor $K^\bullet(X)$ in the variable ringed space $X$, defines a homomorphism of rings
\begin{equation}\tag{5.1}
K^\bullet(X)\longrightarrow K^\bullet(X^{\mathrm{top}}).
\end{equation}
For simplicity suppose that $X^{\mathrm{top}}$ is compact; then, as in the case of a scheme admitting an ample invertible module, one proves (Expos\'e II) that every perfect complex on $X^{\mathrm{top}}$ is globally isomorphic to a bounded complex, locally free in each degree, i.e. to a complex of complex vector bundles on $X^{\mathrm{top}}$, and hence that $K^\bullet(X^{\mathrm{top}})$ is canonically isomorphic to the well-known group $K(X^{\mathrm{top}})$ of topologists [1], defined in terms of complex vector bundles on the compact space $X^{\mathrm{top}}$. Thus one may interpret (5.1) as a homomorphism from the $K^\bullet(X)$ defined in terms of the analytic structure to the well-known ring $K(X^{\mathrm{top}})$ of topologists. If now
\[
f:X\longrightarrow Y
\]
is a morphism of compact nonsingular analytic spaces, one knows how to associate to it, thanks to Atiyah--Hirzebruch [1], a group homomorphism
\begin{equation}\tag{5.2}
f^{\mathrm{top}}_*:K^\bullet(X^{\mathrm{top}})\longrightarrow K^\bullet(Y^{\mathrm{top}}).
\end{equation}
Using likewise the homomorphism $f_*$ of (3.1), defined by formula (3.2), which here implicitly uses Grauert's finiteness theorem [4] for a proper morphism of analytic spaces, one obtains a square of natural homomorphisms involving only groups of type $K^\bullet$, to the exclusion of rings of Chow-ring or cohomology type:
\begin{equation}\tag{5.3}
\begin{array}{ccc}
K^\bullet(X) & \longrightarrow & K^\bullet(X^{\mathrm{top}}) \\
\downarrow f_* && \downarrow f^{\mathrm{top}}_* \\
K^\bullet(Y) & \longrightarrow & K^\bullet(Y^{\mathrm{top}}).
\end{array}
\end{equation}
The Riemann--Roch type formula to which we are alluding consists in asserting the commutativity of this square. Note the difference in nature between this assertion, which does not neglect torsion phenomena and is situated in a group $K^\bullet(Y^{\mathrm{top}})$ of an essentially discrete nature, and the Riemann--Roch formula of the type considered in no. 4, which is situated in a group of ``Chow'' type no longer of a discrete nature, but which, on the other hand, must be tensored with $\mathbb Q$. Let us note that formula (5.3) is proved at any rate when $X$ and $Y$ are nonsingular projective varieties [7], as is the formula of the type considered in no. 4, proved in this case in [RRR]. What is involved here is a plausible variant of the Atiyah--Singer index theorem, and of its generalizations due to Shih [7] and Atiyah, which obviously calls for a common generalization. Once the commutativity of (5.3) is proved, one would immediately deduce from it, using the ``differentiable Riemann--Roch theorem'' of Atiyah--Hirzebruch [1], a cohomological variant of the complex analytic Riemann--Roch formula, namely the commutativity of the diagram
\begin{equation}\tag{5.4}
\begin{array}{ccc}
K^\bullet(X) & \xrightarrow{\;x\mapsto \ch(x)\Todd(T_X)\;} & H^{2\bullet}(X,\mathbb Q) \\
\downarrow f_* && \downarrow f_* \\
K^\bullet(Y) & \xrightarrow{\;y\mapsto \ch(y)\Todd(T_Y)\;} & H^{2\bullet}(Y,\mathbb Q),
\end{array}
\end{equation}
where the horizontal arrows are deduced from cohomological Chern classes, themselves deduced from the homomorphism (5.1) and from the ordinary Chern-class homomorphism
\[
K^\bullet(X^{\mathrm{top}})\longrightarrow H^{2\bullet}(X,\mathbb Z),
\]
the first vertical arrow is the same as in (5.3), and the second is the Gysin homomorphism in rational cohomology.

\subsection*{6. The covariant invariant $K_\bullet(X)$}
In all that precedes, scarcely anything has been said except about the contravariant invariant $K^\bullet(X)$, to the exclusion of the covariant invariant $K_\bullet(X)$ whose existence was indicated in 3.1. If one regards $K^\bullet(X)$ as giving an analogue of the integral cohomology ring of a topological space, one should regard $K_\bullet(X)$ as giving an analogue of integral homology. The module structure of $K_\bullet(X)$ over $K^\bullet(X)$ then corresponds to the cap-product operation. The fact that, when $X$ is regular, the natural homomorphism $K^\bullet(X)\to K_\bullet(X)$ is an isomorphism should be regarded as the analogue of the Poincar\'e duality theorem for oriented topological varieties, asserting that cap product with the fundamental homology class defines an isomorphism from cohomology to homology. These analogies moreover suggest the usefulness of also defining, for a finite polyhedron $X$, say, a covariant invariant $K_\bullet(X)$ whose relations to integral homology -- spectral sequence, Chern homomorphism, and so on -- would be analogous to those existing between $K^\bullet(X)$ and integral cohomology. Nor does it seem excluded that there might exist, in terms of these groups $K_\bullet(X)$, variants of the differentiable Riemann--Roch theorem valid for spaces more general than compact differentiable manifolds.

Returning to the case where $X$ is an arbitrary noetherian scheme, one should endow $K_\bullet(X)$ with an increasing filtration, $\Fil_iK_\bullet(X)$ being generated by classes $\cl_\bullet(F)$, with $F$ a coherent sheaf on $X$ such that $\dim \operatorname{Supp}F\le i$. One will prove (Expos\'e X) the compatibility of this filtration with the module structure and with the filtration of $K^\bullet(X)$, i.e. the formula
\begin{equation}\tag{6.1}
\Fil^iK^\bullet(X)\cdot \Fil_jK_\bullet(X)\subset \Fil_{j-i}K_\bullet(X),
\end{equation}
which allows the graded group associated with the filtered module $K_\bullet(X)$, denoted $\Gr_\bullet(X)$, to be endowed with a module structure over $\Gr^\bullet(X)$. Its elements may in fact be interpreted as classes of algebraic cycles on $X$, for an equivalence relation less fine than rational equivalence when $X$ is of finite type over a base field $k$. If
\[
f:X\longrightarrow Y
\]
is a proper morphism of noetherian schemes, then
\begin{equation}\tag{6.2}
f_*:K_\bullet(X)\longrightarrow K_\bullet(Y)
\end{equation}
is compatible with the filtrations and defines a homomorphism of graded groups
\begin{equation}\tag{6.3}
f_*:\Gr_\bullet(X)\longrightarrow \Gr_\bullet(Y),
\end{equation}
giving rise to a ``projection formula'', i.e. a formula which says that it is $\Gr^\bullet(Y)$-linear; this formula is obtained by passage to associated gradeds from the analogous projection formula for the homomorphism (6.2), which is $K^\bullet(Y)$-linear.

From the point of view proposed in the seminar, a manageable theory of intersections on noetherian schemes $X$, not necessarily regular, should consist in the combined study of the contravariant invariants $K^\bullet(X),\Gr^\bullet(X)$ and the covariant invariants $K_\bullet(X),\Gr_\bullet(X)$. These play roles of essentially different nature, and their properties complement each other. Thus the covariant groups $K_\bullet(X),\Gr_\bullet(X)$ are easier to compute for certain nonproper fibrations, thanks to the ``homotopy exact sequence'' which one can establish for them ([RRR] and Expos\'e IX). For example, one proves (Expos\'e IX) canonical isomorphisms such as
\[
K_\bullet(X[t])\simeq K_\bullet(X),\qquad \Gr_\bullet(X[t])\simeq \Gr_\bullet(X),
\]
which fail in general for $K^\bullet$ when $X$ is a nonregular scheme. On the other hand, $K^\bullet(X)$ and $\Gr^\bullet(X)$ lend themselves well to computations, thanks to their ring structures.

When one considers schemes $X$ proper over a field, the projection
\[
\pi_X:X\longrightarrow (\operatorname{point})
\]
defines a homomorphism, the ``Euler--Poincar\'e characteristic'',
\begin{equation}\tag{6.4}
\chi_X=\pi_{X*}:K_\bullet(X)\longrightarrow K_\bullet(\operatorname{point})=\mathbb Z,
\end{equation}
which induces on the subgroup $\Gr_0(X)=\Fil_0K_\bullet(X)$ the homomorphism ``degree of zero-cycles''
\begin{equation}\tag{6.5}
\deg_X:\Gr_0(X)\longrightarrow \mathbb Z.
\end{equation}
The multiplication homomorphism
\[
\Gr^i(X)\times \Gr_i(X)\longrightarrow \Gr_0(X),
\]
composed with the degree homomorphism (6.5), then furnishes a pairing
\begin{equation}\tag{6.6}
\Gr^i(X)\times \Gr_i(X)\longrightarrow \mathbb Z,
\end{equation}
which explains that, in certain respects, $\Gr^\bullet(X)$ and $\Gr_\bullet(X)$ play dual roles, just as cohomology and homology do. Thus, if $f:X\to Y$ is a morphism of proper algebraic schemes, the two homomorphisms
\[
f^*:\Gr^\bullet(Y)\longrightarrow \Gr^\bullet(X),\qquad f_*:\Gr_\bullet(X)\longrightarrow \Gr_\bullet(Y)
\]
are formally transposes of one another, in the sense of the pairings (6.6).

The homomorphisms (6.4) and (6.5), and the pairing (6.6) deduced from them, may be regarded as the source of numerical relations in intersection theory on proper algebraic schemes over a given field $k$. This point of view will be developed in some detail in Expos\'e X, which is essentially independent of Expos\'es VI, VII, VIII, centred on the proof of the Riemann--Roch theorem, and of Expos\'e IX. Some notions concerning the behaviour of $K^\bullet(X)$ and $\Gr^\bullet(X)$ under specialization will also be given there.

\subsection*{7. Picard groups and the determinant functor}
As in every intersection theory, it is important to specify the place occupied in the theory by the Picard group, or group of classes of divisors in the classical terminology. Subject to a slight restriction, automatically satisfied when $X$ is for instance quasi-compact, one establishes in Expos\'e X a canonical isomorphism
\begin{equation}\tag{7.1}
c_1:\Pic(X)\xrightarrow{\sim}\Gr^1(X),
\end{equation}
which shows the exact role of the Picard group in the theory we are proposing. This role justifies a more detailed study of the Picard group in Expos\'es XII and XIII, which are logically independent of the text of the seminar. Kleiman, in Expos\'e XIII, studies numerical equivalence and certain finiteness theorems for the Picard group, announced without proof in [6], using purely numerical techniques, not relying on the theory of the Picard functor and its representability, due to himself and Mumford, and markedly simpler than Grothendieck's initial proofs. The only result refractory to Kleiman's methods, and which he is obliged to use, is [6, C-07 (i)]; its proof and certain refinements, including certain representability theorems for the Picard functor, occupy Expos\'e XII. The latter is moreover of a spirit and technique entirely foreign to the rest of the seminar, and is connected with [5]; this expos\'e is logically independent of all the others.

Finally, in Expos\'e XI, of a nature appreciably more general than XII and XIII and closer to the general spirit of the seminar, one studies, on an arbitrary locally ringed topos, a remarkable functor called the ``determinant functor'':
\begin{equation}\tag{7.2}
\det^\bullet:\Parf_{\mathrm{is}}(X)\longrightarrow \Inv(X)
\end{equation}
from the category of perfect complexes on $X$, whose morphisms are only the isomorphisms in the sense of the derived category $D(X)$, to the category of invertible modules on $X$. This functor is described, up to very little, by the requirement that it be of ``local nature'' and that it be given, for a bounded complex $L$ with locally free components, by the formula
\begin{equation}\tag{7.3}
\det^\bullet(L)=\bigotimes_i \det(L_i)^{(-1)^i},
\end{equation}
where $\det(L_i)$ denotes the exterior power of maximal degree of the locally free module $L_i$. Although this expos\'e is also logically independent of the rest of the seminar, except for Expos\'e I, it nevertheless constitutes an important element of the general ``yoga'' proposed in the seminar and developed essentially in Expos\'es I to XI.

\section*{Bibliography}
\begin{enumerate}[label={[\arabic*]},leftmargin=2.5em]
\item M. Atiyah and F. Hirzebruch, \emph{Vector bundles and homogeneous spaces}, Symposia in Pure Mathematics, vol. 3, Differential Geometry, pp. 7--38, 1961.
\item A. Borel and J.-P. Serre, \emph{Le th\'eor\`eme de Riemann--Roch}, Bull. Soc. math. France, t. 86, pp. 97--136 (1958).
\item H. Grauert, \emph{Ein Theorem der analytischen Garbentheorie}, Pub. Math. no. 5, pp. 5--64 (1960).
\item A. Grothendieck, \emph{Classes de faisceaux et th\'eor\`eme de Riemann--Roch}, mimeographed notes, Princeton 1957, cited as [RRR] and reproduced as an appendix to Expos\'e 0 of the seminar.
\item A. Grothendieck, \emph{Technique de descente et th\'eor\`emes d'existence en G\'eom\'etrie Alg\'ebrique}, in \emph{Fondements de la G\'eom\'etrie Alg\'ebrique}, extracts from the Bourbaki Seminar 1957--1962, Secr\'etariat Math\'ematique, Paris.
\item W. Shih, S\'eminaire de Topologie \`a l'I.H.E.S., 1966/67.
\item M. Atiyah and F. Hirzebruch, \emph{The Riemann--Roch theorem for analytic embeddings}, Topology, vol. 1, pp. 151--166 (1962).
\end{enumerate}
\end{document}
